\section{Delay-moment range and the Fredholm functional}
\label{sec:fredholm-sensitivity}

The two algebraic ingredients are the range of the first delay moment and
the inverse of the network-transverse generator.  We prove their uniform
bounds before identifying the collective cokernel and the resulting
Fredholm functional.  Section~\ref{sec:fold-passage} then shows that this
functional is the derivative appearing in the local RFDE matching problem.

\subsection{The stationary projection and the first delay moment}

For this subsection, fix only a strictly positive probability vector
\(\pi_N\), the space \(E_N=\ker\pi_N^T\), scalars \(\alpha>0\) and
\(K\ne0\), and the delay values.  Put
\(P_{\perp,N}=\Id-\mathbf1\pi_N^T\), and use the norm and constrained
perturbation space \(\mathcal T_N\) from
\cref{eq:rw-network-norm,eq:rw-structural-parameter-space}.
No property of \(P_N\) or its Dobrushin
coefficient is used in the first-moment result.  When differentiating with
respect to the delayed-coupling matrices, write a tangent direction as
\(R=(R_0,\ldots,R_m)\in\mathcal T_N\).  Thus the \(R_k\) below
are tangent matrices satisfying the same constraints as the components
\(\Delta B_k\) in \cref{eq:rw-structural-parameter-space}; they are not a second
perturbation space.  The stationary projection identity is already
\cref{eq:rw-projection-identity}.

Along the affine singular fold history, translation by \(\theta_k\) produces
the factor \(\theta_k/(2\alpha)\).  The resulting term in the
network-transverse equation is the first delay moment

\begin{equation}
 \mathsf S_NR
 =\frac{K}{2\alpha}P_{\perp,N}
       \left(\sum_{k=0}^m\theta_kR_k\right)\mathbf 1
 \in E_N.
 \label{eq:fredholm-delay-moment}
\end{equation}

We first prove the range theorem and its sharp norm bounds.

\begin{proof}[Proof of \cref{thm:rw-first-moment-map}]
Let \(k_-=0\), \(k_+=m\), and
\(\Delta_\theta=\theta_m-\theta_0\).

For \(R\) satisfying \(\pi_N^TR=0\), set

\[
 \mathsf v_N(R)=P_{\perp,N}R\mathbf 1=R\mathbf 1,
 \qquad
 \mathsf j_N(y)=y\pi_N^T.
\]

If \(y\in E_N\), then

\begin{equation}
 \norm{y\pi_N^T}_{\cL(\R^N,\norm{\cdot}_N)}=\norm{y}_N.
 \label{eq:fredholm-rank-one-norm}
\end{equation}

Indeed, the upper bound follows from
\(\abs{\pi_N^Tx}\le\norm{x}_N\), and equality is attained at
\(x=\mathbf1\).  Consequently
\(\norm{\mathsf v_N}=\norm{\mathsf j_N}=1\), and
\(\mathsf v_N\mathsf j_N=\Id_{E_N}\).

Let \(\bar\theta=(\theta_0+\theta_m)/2\).  For
\(R\in\mathcal T_N\),

\[
 \mathsf S_NR
 =\frac{K}{2\alpha}\mathsf v_N
       \left(\sum_k(\theta_k-\bar\theta)R_k\right),
\]

and hence

\[
 \norm{\mathsf S_NR}_N
 \le\frac{|K|\Delta_\theta}{4\alpha}
      \norm{R}_{\mathcal T_N}.
\]

The matrices in the tuple \cref{eq:rw-source-right-inverse} sum to zero,
and the tuple
satisfies \(\mathsf S_N\mathsf Q_Ny=y\).  Its norm is the second quantity
in \cref{eq:rw-source-sharp-norms}.  Since a right inverse satisfies
\(1\le\norm{\mathsf S_N}\norm{\mathsf Q_N}\), the preceding upper bound
is sharp as well.  If all \(\theta_k=\theta_*\), then

\[
 \sum_k\theta_kR_k=\theta_*\sum_kR_k=0,
\]

which proves the last statement.
\end{proof}

The full-range right inverse above answers the unconstrained matrix problem.
For a fixed topology or another linear matrix restriction, the evaluation
map \(R\mapsto R\mathbf1\) gives the range criterion.

\begin{proof}[Proof of \cref{thm:rw-structured-delay-range}]

For the displayed tuple,
\[
 \sum_k\theta_kR_k
 =-(\theta_{k_+}-\theta_{k_-})R.
\]
Since \(\pi_N^TR=0\), one has \(R\mathbf1\in E_N\), and
\cref{eq:rw-structured-range} follows from
\cref{eq:fredholm-delay-moment}.  This also proves the range criterion and
the row-sum statement.  For
\(R=(P_N-\Id)\diag(d)\),
\[
 \pi_N^TR=0,\qquad R\mathbf1=(P_N-\Id)d.
\]
Irreducibility makes the eigenvalue \(1\) of \(P_N\) simple, so
\((P_N-\Id)|_{E_N}\) is invertible.  The support and local nonnegativity
claims follow directly from the form of \(R\).

Under the Dobrushin hypothesis,
\[
 \norm{d_y}_N\le\gamma^{-1}\norm{y}_N.
\]
For \(d\in E_N\) and arbitrary \(x\),
\[
 \osc(d\circ x)\le2\norm{d}_N\norm{x}_N.
\]
Together with
\(\osc\{(P_N-\Id)u\}\le(2-\gamma)\osc(u)\), this gives
\cref{eq:rw-generator-right-inverse-bound}.  Multiplying
\(\mathsf J_N\) by \(-2\alpha/
\{K(\theta_{k_+}-\theta_{k_-})\}\) in layer \(k_-\), and by the opposite
coefficient in layer \(k_+\), gives
\cref{eq:rw-structured-source-inverse-bound}.  Finally, every off-diagonal
entry of the two base layers is \(P_{ij}/2\), while the corresponding entry of
\(R\) is \(P_{ij}d_j\); on the diagonal the base entry is
\((1+P_{jj})/2\) and the perturbation is
\((P_{jj}-1)d_j\).  The condition
\(|\zeta|\norm d_\infty<1/2\) therefore proves the last assertion.
\end{proof}

\subsection{The transverse inverse and the collective cokernel}

The Dobrushin gap supplies the inverse needed to propagate
\cref{eq:fredholm-delay-moment}.  This is an inverse of the transverse
linear block, not an inverse RFDE semiflow.

\begin{lemma}[Uniform inverse on the stationary-mean-zero subspace]
\label[lemma]{lem:fredholm-dobrushin-inverse}
Suppose

\[
 \tau(P_N)\le1-\gamma,
 \qquad \gamma>0,
\]

with \(\gamma\) independent of \(N\).  Then \(A_N\) is invertible on
\(E_N\), and for \(t\ge0\) and \(z\in E_N\),

\begin{align}
 \norm{e^{A_Nt}z}_N
 &\le e^{-D\gamma t}\norm z_N,
 \label{eq:fredholm-dobrushin-semigroup}\\
 \norm{A_N^{-1}}_{E_N\to E_N}
 &\le\frac1{D\gamma},
 &
 \norm{A_Nz}_N
 &\le D(2-\gamma)\norm z_N.
 \label{eq:fredholm-dobrushin-operator-bounds}
\end{align}
\end{lemma}

\begin{proof}
On \(E_N\), the norm is the oscillation seminorm, and Dobrushin
contraction gives

\[
 \norm{P_N^jz}_N\le(1-\gamma)^j\norm z_N.
\]

The Poisson expansion therefore yields

\[
 \norm{e^{A_Nt}z}_N
 \le e^{-Dt}\sum_{j=0}^\infty
       \frac{(Dt)^j}{j!}(1-\gamma)^j\norm z_N
 =e^{-D\gamma t}\norm z_N.
\]

Thus

\[
 A_N^{-1}=-\int_0^\infty e^{A_Nt}\,\dd t
\]

and the inverse bound follows.  Finally,

\[
 \norm{A_Nz}_N
 \le D\{\norm{P_Nz}_N+\norm z_N\}
 \le D(2-\gamma)\norm z_N.
\]
\end{proof}

\begin{proof}[Proof of \cref{prop:rw-growing-families}]

For the mixing family,
\[
 \osc(P_N^{\rm mix}x)=(1-\rho)\osc(x),
 \qquad
 A_N=-D\rho\Id\quad\hbox{on }E_N.
\]
Thus \(\tau(P_N^{\rm mix})=1-\rho\).  Direct summation gives
\[
 \pi_N^Tq_N=0,
 \qquad \pi_N^Tq_N^{\circ2}=1,
 \qquad \max_i|q_{i,N}|<\sqrt3.
\]
For \((\theta_0,\theta_1)=(1,2)\) and
\((R_0,R_1)=(q_N\pi_N^T,-q_N\pi_N^T)\),
\[
 P_{\perp,N}\left(\theta_0R_0+\theta_1R_1\right)\mathbf1=-q_N.
\]
Since \(\alpha=\pi_N^Tc_N=1\), insertion in
\cref{eq:fredholm-curvature-lambda} yields
\[
 \Lambda_N(R_0,R_1)
 =-\frac{K}{2D\rho}\pi_N^T\diag(\mathbf1+\sigma q_N)q_N
 =-\frac{K\sigma}{2D\rho}.
\]
The strict bound on \(q_N\) proves the common positivity of \(c_N\), and its
entries are distinct by construction.  By
\cref{eq:fredholm-rank-one-norm},
\[
 \norm{(R_0,R_1)}_{\mathcal T_N}
 =2\norm{q_N\pi_N^T}=2\osc(q_N)<4\sqrt3,
\]
which proves the uniform bound on the perturbation direction.

For the cyclic family, two rows with disjoint supports give
\(\tau(P_N^{\rm cyc})=1\) when \(N\ge4\).  Let \(N=2n\), and take
\(y_i=1\) for \(1\le i\le n\) and \(y_i=-1\) otherwise.  Then
\(y\in E_N\) and \(\norm{y}_N=2\).  With the convention
\((C_Nz)_i=z_{i+1}\) and cyclic indices, solving
\[
 D\rho(C_N-\Id)z=y
\]
gives, up to an additive constant, a sequence whose oscillation is
\(n/(D\rho)\).  Consequently the inverse norm is at least
\(n/(2D\rho)=N/(4D\rho)\), as claimed.
\end{proof}

For the collective fold equation, the relevant whole-line operator is

\[
 L_0(U,V)=(U'-sU-V,V'+U).
\]

Fix \(c>0\) and \(M\ge0\), and use the weighted source space

\[
 \mathcal Y_0
 =\left\{f\in C(\R;\R^2):
   \sup_{s\in\R}e^{-c|s|}\langle s\rangle^{-M}|f(s)|<\infty\right\}.
\]

Let \(\mathcal D_0\) be the phase-fixed graph domain consisting of
\(\xi=(U,V)\in C^1(\R;\R^2)\) with \(U(0)=0\), \(L_0\xi\in\mathcal Y_0\),
and

\[
 \sup_{s\in\R}e^{-c|s|}\langle s\rangle^{-M-1}|\xi(s)|<\infty.
\]

Equip \(\mathcal D_0\) with the corresponding graph norm.

\begin{lemma}[Cokernel of the phase-fixed fold operator]
\label[lemma]{lem:fredholm-collective-cokernel}
The operator \(L_0:\mathcal D_0\to\mathcal Y_0\) is injective, has closed
range of codimension one, and

\begin{equation}
 \Ran L_0=\Ker\ell_0,
 \qquad
 \ell_0f
 =\int_{\R}e^{-s^2/2}\{sf_1(s)+f_2(s)\}\,\dd s.
 \label{eq:fredholm-gaussian-cokernel}
\end{equation}

It has a bounded inverse from \(\Ker\ell_0\) to \(\mathcal D_0\).
\end{lemma}

\begin{proof}
Put

\[
 E(s)=e^{s^2/2},\qquad I(s)=\int_0^sE(t)\,\dd t,
\]

and introduce the fundamental solutions

\[
 \tau(s)=(-1,s)^T,
 \qquad
 n(s)=(I(s),E(s)-sI(s))^T.
\]

Writing \(\xi=a\tau+bn\), variation of constants gives

\begin{equation}
 a'=E^{-1}\{-(E-sI)f_1+If_2\},
 \qquad
 b'=E^{-1}(sf_1+f_2).
 \label{eq:fredholm-fold-variation}
\end{equation}

Membership in \(\mathcal D_0\) forces
\(b(s)=E(s)^{-1}\{sU(s)+V(s)\}\to0\) as \(s\to\pm\infty\).
Integration of the second equation therefore gives
\(\ell_0(L_0\xi)=0\).  Conversely, if \(\ell_0f=0\), choose

\[
 a(s)=\int_0^sE(t)^{-1}
       \{-(E(t)-tI(t))f_1(t)+I(t)f_2(t)\}\,\dd t,
\]

and

\[
 b(s)=\int_{-\infty}^sE(t)^{-1}
       \{tf_1(t)+f_2(t)\}\,\dd t
     =-\int_s^\infty E(t)^{-1}
       \{tf_1(t)+f_2(t)\}\,\dd t.
\]

The standard Gaussian tail estimates, together with

\[
 |I(s)|\le CE(s)\langle s\rangle^{-1},
 \qquad
 |E(s)-sI(s)|\le CE(s),
\]

give a bounded map \(f\mapsto a\tau+bn\) from
\(\Ker\ell_0\) to \(\mathcal D_0\).  If \(f=0\), the weighted domain
excludes the Gaussian mode \(n\), while \(U(0)=0\) removes the tangent mode
\(\tau\).  Hence \(L_0\) is injective.  The bounded inverse on
\(\Ker\ell_0\), and the fact that \(\ell_0\) is a nonzero bounded
functional, prove the closed-range and codimension statements.
\end{proof}

We now state the Lyapunov--Schmidt reduction in the form used below.  It is
important that the projection identity by itself does not remove possible
derivatives of the phase, endpoint, or matching conditions with respect to
the delayed-coupling matrices.  Such terms
must either vanish or belong to the collective range.

\begin{theorem}[Lyapunov--Schmidt reduction for a triangular block system]
\label{thm:fredholm-reduction}
For each \(N\), let

\[
 L_{\perp,N}:\mathcal D_{\perp,N}\to\mathcal Y_{\perp,N},
 \qquad
 L_{\parallel,N}:\mathcal D_{\parallel,N}\to\mathcal Y_{\parallel,N}
\]

be bounded linear operators between real Banach spaces.  Let

\[
 \mathsf S_N:\mathcal T_N\to\mathcal Y_{\perp,N},
 \qquad
 \mathcal B_N:\mathcal D_{\perp,N}\to\mathcal Y_{\parallel,N},
 \qquad
 d_N:\mathcal T_N\to\mathcal Y_{\parallel,N}
\]

be bounded, and let \(f_{\nu,N}\in\mathcal Y_{\parallel,N}\).  Assume that

\begin{enumerate}[label=\textup{(F\arabic*)}]
 \item \(L_{\perp,N}\) is a bounded isomorphism;
 \item \(\Ran L_{\parallel,N}\) is closed of codimension one;
 \item \(f_{\nu,N}\notin\Ran L_{\parallel,N}\);
 \item \(d_NR\in\Ran L_{\parallel,N}\) for every \(R\in\mathcal T_N\).
\end{enumerate}

For the first-variation equations

\begin{equation}
 L_{\perp,N}h=\mathsf S_NR,
 \qquad
 L_{\parallel,N}x+f_{\nu,N}\lambda
      +\mathcal B_Nh+d_NR=0,
 \label{eq:fredholm-block-system}
\end{equation}

there is a unique scalar-valued sensitivity functional
\(\Lambda_N\in\mathcal T_N^*\) such
that the solvability condition is \(\lambda=\Lambda_N(R)\).  If
\(\ell_N\in\mathcal Y_{\parallel,N}^*\setminus\{0\}\) satisfies
\(\Ker\ell_N=\Ran L_{\parallel,N}\), then

\begin{equation}
 \Lambda_N(R)
 =-\frac{\ell_N\mathcal B_N
                 L_{\perp,N}^{-1}\mathsf S_NR}
          {\ell_Nf_{\nu,N}}.
 \label{eq:fredholm-lambda}
\end{equation}
For a fixed direction \(R\), the scalar \(\Lambda_N(R)\) is the corresponding
first-order coefficient.

The quotient is independent of the nonzero normalization of \(\ell_N\), of
the range-valued term \(d_NR\), and of adding any further element of
\(\Ran L_{\parallel,N}\) to the collective right-hand side.  Moreover,

\begin{equation}
 \norm{\Lambda_N}
 \le
 \frac{\norm{\mathcal B_N}
       \norm{L_{\perp,N}^{-1}}
       \norm{\mathsf S_N}}
 {\operatorname{dist}
    (f_{\nu,N},\Ran L_{\parallel,N})}.
 \label{eq:fredholm-lambda-bound}
\end{equation}
\end{theorem}

\begin{proof}
The first line of \eqref{eq:fredholm-block-system} gives
\(h=L_{\perp,N}^{-1}\mathsf S_NR\).  Applying \(\ell_N\) to the second
equation kills \(L_{\parallel,N}x\) and \(d_NR\), and gives

\[
 (\ell_Nf_{\nu,N})\lambda
 +\ell_N\mathcal B_NL_{\perp,N}^{-1}\mathsf S_NR=0.
\]

Since the cokernel is one dimensional and
\([f_{\nu,N}]\ne0\), this determines \(\lambda\) uniquely and proves
\cref{eq:fredholm-lambda}.  The quotient-norm estimate in the cokernel is
\cref{eq:fredholm-lambda-bound}.  The remaining invariance statements
follow because \(\ell_N\) annihilates the collective range.
\end{proof}

\subsection{The induced functional on the transverse subspace}

The sensitivity functional determines only the linear functional on the
transverse subspace that has entered the one-dimensional Fredholm equation.
It does not determine
the coupling matrix \(P_N\) or an arbitrary history state.

\begin{proposition}[Reconstruction of the transverse linear functional]
\label[proposition]{prop:fredholm-dual-recovery}
Under \cref{thm:fredholm-reduction}, define

\begin{equation}
 \mathfrak r_N(h)
 =-\frac{\ell_N\mathcal B_Nh}{\ell_Nf_{\nu,N}},
 \qquad h\in\mathcal D_{\perp,N}.
 \label{eq:fredholm-return-functional}
\end{equation}

Then

\begin{equation}
 \Lambda_N
 =\mathfrak r_NL_{\perp,N}^{-1}\mathsf S_N.
 \label{eq:fredholm-return-factorization}
\end{equation}

If \(\mathsf S_N\) is onto and has a bounded right inverse
\(\mathsf Q_N:\mathcal Y_{\perp,N}\to\mathcal T_N\), then

\begin{equation}
 \mathfrak r_N(h)
 =\Lambda_N(\mathsf Q_NL_{\perp,N}h),
 \qquad h\in\mathcal D_{\perp,N},
 \label{eq:fredholm-return-reconstruction}
\end{equation}

and

\begin{equation}
 \frac{\norm{\mathfrak r_N}}
      {\norm{\mathsf Q_NL_{\perp,N}}}
 \le\norm{\Lambda_N}
 \le\norm{\mathfrak r_N}
      \norm{L_{\perp,N}^{-1}\mathsf S_N}.
 \label{eq:fredholm-return-condition}
\end{equation}

For the transverse block \(L_{\perp,N}=A_N\) associated with the
row-stochastic coupling, use \cref{eq:rw-source-right-inverse} and write
\(\Delta_\theta=\theta_m-\theta_0\).  Then

\begin{equation}
 \frac{|K|\Delta_\theta}{4\alpha D(2-\gamma)}
       \norm{\mathfrak r_N}
 \le\norm{\Lambda_N}
 \le
 \frac{|K|\Delta_\theta}{4\alpha D\gamma}
       \norm{\mathfrak r_N}.
 \label{eq:fredholm-markov-dual-condition}
\end{equation}
\end{proposition}

\begin{proof}
Equation~\eqref{eq:fredholm-return-factorization} is
\cref{eq:fredholm-lambda} rewritten using
\cref{eq:fredholm-return-functional}.  If
\(\mathsf S_N\mathsf Q_N=\Id\), then

\[
 \Lambda_N\mathsf Q_NL_{\perp,N}
 =\mathfrak r_NL_{\perp,N}^{-1}
    \mathsf S_N\mathsf Q_NL_{\perp,N}
 =\mathfrak r_N,
\]

which proves \cref{eq:fredholm-return-reconstruction}.  The two norm
bounds follow from this identity and
\cref{eq:fredholm-return-factorization}.  Finally,
\cref{eq:rw-source-sharp-norms,eq:fredholm-dobrushin-operator-bounds}
give

\[
 \norm{A_N^{-1}\mathsf S_N}
 \le\frac{|K|\Delta_\theta}{4\alpha D\gamma},
 \qquad
 \norm{\mathsf Q_NA_N}
 \le\frac{4\alpha D(2-\gamma)}{|K|\Delta_\theta},
\]

and hence \cref{eq:fredholm-markov-dual-condition}.
\end{proof}

For the polynomial fold equation in \cref{sec:rw-model-results},
the collective forcing produced by a constant
transverse displacement \(h\in E_N\) is

\begin{equation}
 (\mathcal B_Nh)(s)
 =\left(\frac{s}{\alpha}\pi_N^T\diag(c_N)h,\,0\right)^T,
 \qquad
 f_{\nu,N}(s)=(0,1)^T.
 \label{eq:fredholm-curvature-profiles}
\end{equation}

Since

\[
 \ell_0f_{\nu,N}=\sqrt{2\pi},
 \qquad
 \ell_0\mathcal B_Nh
 =\frac{\sqrt{2\pi}}{\alpha}
      \pi_N^T\diag(c_N)h,
\]

the corresponding linear functional on \(E_N\) is

\begin{equation}
 \mathfrak r_N^{\rm curv}(h)
 =-\frac1\alpha\pi_N^T\diag(c_N)h.
 \label{eq:fredholm-curvature-return}
\end{equation}

\begin{corollary}[Specialization to the polynomial network]
\label[corollary]{cor:fredholm-curvature-specialization}
For the network equation in \cref{eq:rw-voltage-rfde},

\begin{equation}
 \Lambda_N(R)
 =-\frac{K}{2\alpha^2}
   \pi_N^T\diag(c_N)A_N^{-1}P_{\perp,N}
   \left(\sum_{k=0}^m\theta_kR_k\right)\mathbf1.
 \label{eq:fredholm-curvature-lambda}
\end{equation}

For two chosen delays \(\theta_{k_-}<\theta_{k_+}\), write
\(\Delta_\theta=\theta_{k_+}-\theta_{k_-}\).
If \(P_N\) is irreducible and \(d\in E_N\), let
\(R(d)\in\mathcal T_N\) be the tangent direction with

\begin{equation}
 R_{k_-}(d)=(P_N-\Id)\diag(d),
 \qquad
 R_{k_+}(d)=-R_{k_-}(d),
 \label{eq:fredholm-generator-probe}
\end{equation}
and all other components zero.  Then

\begin{equation}
 \Lambda_N(R(d))
 =\frac{K\Delta_\theta}{2\alpha^2D}
   \left\langle c_N-\alpha\mathbf1,d\right\rangle_{\pi_N},
 \qquad
 \langle x,y\rangle_{\pi_N}=\pi_N^T(x\circ y).
 \label{eq:fredholm-curvature-pairing}
\end{equation}

Consequently this restriction is nonzero exactly when
\(c_N\ne\alpha\mathbf1\).

\end{corollary}

\begin{proof}
Combining
\cref{eq:fredholm-delay-moment,eq:fredholm-curvature-return,eq:fredholm-return-factorization}
gives
\cref{eq:fredholm-curvature-lambda}.  For
\cref{eq:fredholm-generator-probe},

\[
 \sum_k\theta_kR_k(d)
 =-\Delta_\theta(P_N-\Id)\diag(d).
\]

Since \(d\in E_N\),

\[
 A_N^{-1}(P_N-\Id)d=D^{-1}d,
\]

and \cref{eq:fredholm-curvature-pairing} follows.  The weighted pairing is
nondegenerate on \(E_N\); taking \(d=c_N-\alpha\mathbf1\) gives the positive
stationary variance whenever the quadratic coefficients are heterogeneous.
\end{proof}
